%---------------------------------------------------------------------------%
%-                                                                         -%
%-                           中文摘要                                      -%
%-                                                                         -%
%---------------------------------------------------------------------------%
% 摘要是论文的重要组成部分，是对论文研究内容和成果的高度概括
% 摘要应包括以下几个方面：
% 1. 研究目的和背景
% 2. 研究方法和过程
% 3. 研究结果和结论
% 4. 研究的创新点和意义

\begin{chineseabstract}
	{APPT；AI驱动的幻灯片协作编辑平台；实时协作；Yjs；Slidev；权限管理}随着信息技术的快速发展和数字化办公的普及，传统幻灯片制作工具在团队协作、版本控制和智能辅助方面存在明显局限性。为解决这些问题，本研究设计并实现了APPT（AI驱动的幻灯片协作编辑平台），该平台基于现代Web技术栈，整合实时协作技术、权限管理系统、AI辅助功能和简化版本控制，为技术团队和教育工作者提供高效的幻灯片制作与协作解决方案。

平台采用React 19 + TypeScript + Vite构建前端，集成CodeMirror 6编辑器和Yjs协作库；后端采用NestJS框架，提供RESTful API和基于Hocuspocus的WebSocket协作服务。系统实现了实时协作编辑功能，支持多用户同时编辑和实时光标追踪；设计了基于RBAC的四级权限管理系统，实现细粒度的文档访问控制；集成了AI辅助功能，支持基于大语言模型的幻灯片内容生成和语法适配；构建了简化的版本控制系统，提供版本历史查看、差异对比和回滚功能。

APPT平台具有创新性的协作架构设计，通过Yjs CRDT算法实现无冲突的实时协作编辑，降低非技术用户的使用门槛。平台在实际应用中表现出良好的稳定性和用户体验，能够显著提高团队协作效率，降低技术使用门槛，增强内容创作能力，促进知识共享与管理。该研究成果为智能幻灯片协作平台的设计与开发提供了理论参考和实践指导。

\end{chineseabstract}
%---------------------------------------------------------------------------%
%-                                                                         -%
%-                           英文摘要                                      -%
%-                                                                         -%
%---------------------------------------------------------------------------%
% The abstract is an important part of the thesis, providing a high-level summary
% of the research content and results. It should include:
% 1. Research purpose and background
% 2. Research methods and process
% 3. Research results and conclusions
% 4. Innovation points and significance of the research

\begin{englishabstract}
		{APPT; AI-driven Slideshow Collaboration Platform; Real-time Collaboration; Yjs; Slidev; Permission Management}With the rapid development of information technology and the popularization of digital office, traditional slideshow creation tools have significant limitations in team collaboration, version control, and intelligent assistance. To address these issues, this research designs and implements APPT (AI-driven Slideshow Collaboration Platform), a modern web technology-based platform that integrates real-time collaboration technology, permission management systems, AI-assisted features, and simplified version control, providing efficient slideshow creation and collaboration solutions for technical teams and educators.

The platform adopts React 19 + TypeScript + Vite for front-end development, integrating CodeMirror 6 editor and Yjs collaboration library; the back-end uses NestJS framework to provide RESTful API and Hocuspocus-based WebSocket collaboration services. The system implements real-time collaborative editing functionality, supporting multi-user simultaneous editing and real-time cursor tracking; designs an RBAC-based four-level permission management system for fine-grained document access control; integrates AI-assisted features supporting large language model-based slideshow content generation and syntax adaptation; constructs a simplified version control system providing version history viewing, difference comparison, and rollback functionality.

APPT platform features innovative collaboration architecture design, implementing conflict-free real-time collaborative editing through Yjs CRDT algorithms, lowering the barrier for non-technical users. The platform demonstrates good stability and user experience in practical applications, significantly improving team collaboration efficiency, lowering technical barriers, enhancing content creation capabilities, and promoting knowledge sharing and management. This research provides theoretical reference and practical guidance for the design and development of intelligent slideshow collaboration platforms.

\end{englishabstract}